% Source: https://artofproblemsolving.com/community/c3h1466161_practice_from_an_old_mock_not_made_by_me
% Stable identity: aops-topic-1466161-q24
\textbf{The area of a lattice triangle is a multiple of $\dfrac{1}{2}$.} If the vertices are $(x_1,y_1)$, $(x_2,y_2)$, $(x_3,y_3)$ with integer coordinates, the shoelace formula gives
\[
\text{Area}=\dfrac{1}{2}\bigl|x_1(y_2-y_3)+x_2(y_3-y_1)+x_3(y_1-y_2)\bigr|,
\]
and the expression inside the absolute value is an integer. So the area is rational.

\textbf{The triangle with sides $8,10,12$ has irrational area.} Its semiperimeter is $s=15$, and Heron's formula gives
\[
\text{Area}=\sqrt{15(15-8)(15-10)(15-12)}=\sqrt{15\cdot7\cdot5\cdot3}=\sqrt{1575}=15\sqrt{7},
\]
which is irrational.

\begin{asy}
usepackage("amsmath");
size(5.5cm);
defaultpen(fontsize(10pt));
pair A=(0,0),B=(12,0),C=(7.5,sqrt(100-7.5^2));
draw(A--B--C--cycle,linewidth(0.9));
dot(A);dot(B);dot(C);
label("$A$",A,SW);label("$B$",B,SE);label("$C$",C,N);
label("$12$",(A+B)/2,S);label("$10$",(A+C)/2,NW);label("$8$",(B+C)/2,NE);
label("\mbox{}",(min(currentpicture,user=true)+max(currentpicture,user=true))/2,black); // keeps the text after the figure black when LaTeX typesets the labels inline
\end{asy}

Every triangle with side lengths $8$, $10$, $12$ is congruent to this one and has area $15\sqrt{7}$, which is not rational. Hence no such triangle can have all three vertices at lattice points. $\blacksquare$

(So a problem asking for the number of lattice points inside such a triangle has no answer; Pick's theorem cannot be applied because the triangle does not exist on the lattice.)
